\documentclass[../main.tex]{subfiles}
\begin{document}
\setlist[itemize]{topsep=2pt,itemsep=1pt,parsep=0pt,partopsep=0pt}
\setlength{\parskip}{3pt}

\chapter{Report Writing Guide}
\label{ch:reportguide}

The RM report is where the whole course is written into one document. Each section is generated by a specific earlier week; this chapter maps those connections and flags the common pitfalls.

\section{Section-to-Concept Map}
\label{sec:report-map}

\autoref{tab:reportmap} traces each report section to its source week and concept.

\begin{table}[H]
\centering
\small
\begin{tabular}{@{}p{2.7cm}p{4.3cm}p{8.1cm}@{}}
\toprule
\textbf{Report section} & \textbf{Generated by} & \textbf{What it must contain} \\
\midrule
Introduction (head) &
Week~1 problem and research ideas (\autoref{sec:research-process}, \autoref{sec:research-ideas}); value of research (\autoref{sec:eop}) &
Why the topic matters, the literature gap, and its relevance to science or industry, kept proportional to the body. \\
\addlinespace
Research Questions &
Week~2 objective and questions (\autoref{sec:objectives-rqs}) &
The research objective and at least two open-ended sub-questions that are relevant, researchable, anchored, and precise. \\
\addlinespace
Methods, Tools, and Expected Results &
Week~2 technical design (\autoref{sec:strategy-materials}); Week~3 sampling and data (\autoref{sec:sampling}, \autoref{sec:data-sources}); Week~4 V\&V (\autoref{sec:vv}) &
Strategy and materials, sampling and data approach, variables and measurement, and the verification and validation plan. \\
\addlinespace
Planning &
Week~2 planning (\autoref{sec:strategy-materials}); Week~5 project and data management (\autoref{sec:project-mgmt}, \autoref{sec:dmp}) &
WBS decomposed to estimable tasks, Gantt chart with critical path, required resources, and data management approach. \\
\addlinespace
Conclusions (feet) &
Week~6 report structure (\autoref{sec:reports}) &
Concise summary of the entire plan, self-contained, introducing no new information. \\
\bottomrule
\end{tabular}
\caption{Report sections mapped to course weeks and concepts.}
\label{tab:reportmap}
\end{table}

\section{Section Notes and Common Pitfalls}
\label{sec:report-notes}

\textbf{Introduction.} State the literature gap and why closing it matters. Pitfalls: an introduction longer than the body it supports; asserting importance without grounding it in the literature.

\textbf{Research questions.} Present the objective and at least two open-ended sub-questions passing the RRAP test. Pitfalls: importing the client's questions; questions so vague that the methods section cannot be specified.

\textbf{Methods, tools, and expected results.} Connect each method to the variables it produces. Pitfalls: describing methods in isolation from the questions they serve; omitting validation entirely.

\textbf{Planning.} Present a WBS, a Gantt chart with a visible critical path, resource requirements, and the data management approach. Pitfalls: no critical path (hides where delay costs time); no flexibility (will not survive contact with the project).

\textbf{Conclusions.} Summarise the entire plan so completely that a reader could grasp it from the conclusions alone. Pitfall: introducing new information that belongs in the body.

\end{document}
